% ch-a4-01.tex —— 第 4 篇（Emerton–Gee–Hellmann）篇首 + §1，底本 PDF 180–191 = 印刷页 167–178
\chapter{范畴 $p$-进 Langlands 纲领导引：导论}
\ysrc{167}
\begin{center}
{\small\itshape Matthew Emerton \quad Toby Gee \quad Eugen Hellmann}
\end{center}
\medskip
\noindent\textbf{摘要.} 我们介绍 $p$-进 Langlands 纲领的「范畴」进路，兼及「Banach」与「解析」两种设定。
\footnote{2020 年数学主题分类：主分类 11R39。}
\footnote{第一作者部分受 NSF 资助（DMS-1902307、DMS-1952705、DMS-2201242）。第二作者部分受一项 ERC Advanced Grant 资助。本项目获欧洲研究理事会（ERC）在欧盟「地平线 2020」研究与创新计划下的资助（资助协议号 884596）。第一、第二作者均部分受 Simons 基金会「代数、几何与拓扑中的完美化」（Perfection in Algebra, Geometry and Topology）合作项目资助。第三作者受德国「卓越战略」EXC 2044-390685587「Mathematics Münster: Dynamics--Geometry--Structure」以及 DFG 的 CRC 1442「Geometry: Deformations and Rigidity」资助。\textcopyright 2025 American Mathematical Society。}

\medskip
\noindent\textbf{目录.}
\begin{list}{}{\leftmargin=2.6em\labelwidth=2.4em\labelsep=0.4em\itemindent=-2.6em\itemsep=0.1em}
\item 1.\quad 引言
\item 1.1.\quad Langlands 纲领速览
\item 1.2.\quad $p$-进视角
\item 1.3.\quad 互反律的证明
\item 1.4.\quad 非分歧几何 Langlands 与 Fargues--Scholze 猜想
\item 1.5.\quad $\ell$-进与 $p$-进情形的若干差异
\item 1.6.\quad 来自 Taylor--Wiles 拼接的动机
\item 1.7.\quad 「从伽罗瓦到自守」的方向
\item 1.8.\quad 我们略去的话题
\item 1.9.\quad 本讲义导读
\item 1.10.\quad 致谢
\item 2.\quad 记号
\item 2.1.\quad 域
\item 2.2.\quad $p$-进 Hodge 理论
\item 2.3.\quad Hodge--Tate 权与最高权
\item 3.\quad 作为范畴 $p$-进局部 Langlands 之动机的 Taylor--Wiles 拼接
\item 3.1.\quad 拼接的极简介绍
\item 3.2.\quad $p$-进 Langlands：概览与历史
\item 4.\quad $(\varphi,\Gamma)$-模的模叠：Banach 情形
\item 4.1.\quad [EG23] 概述
\item 4.2.\quad 固定行列式变体
\item 5.\quad $(\varphi,\Gamma)$-模的模叠：解析情形
\end{list}
\ysrc{168}
\begin{list}{}{\leftmargin=2.6em\labelwidth=2.4em\labelsep=0.4em\itemindent=-2.6em\itemsep=0.1em}
\item 5.1.\quad Robba 环上的 $(\varphi,\Gamma)$-模
\item 5.2.\quad de Rham 对象的叠
\item 5.3.\quad Drinfeld 型紧化
\item 6.\quad 范畴 $p$-进局部 Langlands 猜想
\item 6.1.\quad 预期：Banach 情形
\item 6.2.\quad 预期：解析情形
\item 7.\quad 范畴 $p$-进 Langlands 的若干已知情形
\item 7.1.\quad $\mathrm{GL}_1$
\item 7.2.\quad $\mathrm{GL}_2(\Q_p)$ 的 Banach 情形——I.\ $\mathcal X$ 的结构
\item 7.3.\quad $\mathrm{GL}_2(\Q_p)$ 的 Banach 情形——II.\ 证明梗概
\item 7.4.\quad $\mathrm{GL}_2(\Q_p)$ 的 Banach 情形——III.\ 例
\item 7.5.\quad 半正交分解
\item 7.6.\quad $\mathrm{GL}_2(\Q_{p^f})$ 的 Banach 情形中的预备
\item 7.7.\quad $\mathrm{GL}_2(\Q_{p^2})$ 的 Banach 情形中的部分结果与例
\item 7.8.\quad 伴随函子
\item 7.9.\quad 整体函数与 Hecke 代数
\item 8.\quad $\ell\neq p$ 情形的范畴局部 Langlands
\item 8.1.\quad C-群
\item 8.2.\quad 局部情形（$\ell\neq p$）的 Langlands 参数
\item 8.3.\quad 范畴局部 Langlands
\item 9.\quad 整体 Langlands 参数的模叠与志村簇的上同调
\item 9.1.\quad 整体模叠
\item 9.2.\quad 整体到局部的映射
\item 9.3.\quad 用范畴 Langlands 计算志村簇的上同调
\item 9.4.\quad 整体情形的一例：$(\mathrm{PGL}_2)/\Q$
\item 9.5.\quad $\mathcal X^{\mathrm{odd}}_{\mathrm{GL}_2,\Q,S}$ 的结构
\item 9.6.\quad 特征簇与超收敛 $p$-进自守形式的层
\item 附录 A.\quad 范畴论背景
\item 附录 B.\quad ind-代数叠上的凝聚层、ind-凝聚层与 pro-凝聚层
\item 附录 C.\quad 局部化与 Morita 理论
\item 附录 D.\quad 拓扑阿贝尔群、环与模
\item 附录 E.\quad $p$-进解析群的表示
\item 附录 F.\quad 群表示的导出模叠
\item 参考文献
\end{list}

\section{引言}

本讲义旨在讨论 Langlands 纲领的若干 $p$-进侧面，尤其是正在浮现的「范畴」视角。我们先作一概述，先讲一些历史，再讲我们的目标；遇到不熟悉的记号，读者可参看第 2 节。

\subsection{Langlands 纲领速览}

Langlands 纲领始于 Langlands 致 Weil 的那封著名信件（例如重印于 [DS15, §6]），其内容在他随后的若干著作中得到展开\ysrc{169}，包括他的文章《自守形式理论中的问题》[Lan70] 以及他关于 Euler 乘积的 Yale 讲义 [Lan71]。该文聚焦于自守表示的函子性猜想及其推论，而 Yale 讲义则解释了他对自守 $L$-函数的构造，并讨论如何证明其预期的解析性质（解析延拓与函数方程）。\footnote{由于 $\mathrm{GL}_d$ 上自守表示的标准 $L$-函数与 Godement--Jacquet [GJ72] 及 Tamagawa [Tam63] 所构造的 $L$-函数一致，而后两者已知具有解析延拓与函数方程，故这些预期性质可由函子性推出。Langlands 的 Yale 讲义讨论了另一条途径，即通过 Eisenstein 级数的常数项。}

在 Yale 讲义中，Langlands 已经触及（尽管只是旁及）如下想法：自守 $L$-函数可能囊括所有动机 $L$-函数。这一想法后来发展成互反律猜想。（该猜想的一种表述见 [Lan79]。）粗略地说，互反律猜想陈述了如下对应：一侧是 $G(\mathbf A_F)$ 的某些自守表示（即代数自守表示），其中 $G$ 是数域 $F$ 上的连通约化代数群；另一侧是 $F$ 上（例如以 $\Q$ 为系数的）动机，其动机伽罗瓦群与 $G$ 的 C-群 ${}^cG$\footnote{C-群由 [BG14] 引入，是 Langlands 所引入的 L-群的加细，更适合处理自守形式与伽罗瓦表示之间的关联问题；下文将回忆其定义。} 密切相关\tn{底本此处作「closely related to to the C-group」，多一 to，疑为笔误；译文按原意处理。}。关于伽罗瓦表示的种种猜想（尤其是 Fontaine--Mazur 猜想 [FM95]）提示：这样的动机又可以等同于 $\ell$-进伽罗瓦表示的相容系统 $\Gal_F\to{}^cG(\overline{\Q}_\ell)$。\footnote{据我们所知，这样的预期只能在 $G=\mathrm{GL}_n$ 的情形精确表述，因为对一般的 $G$，伽罗瓦表示的相容系统之精确定义尚不清楚；例如见 [BHKT19, §6]。}

要精确表述互反律猜想，牵涉许多微妙之处。例如，当 $G$ 不是某个 $\mathrm{GL}_d$ 时，须顾及 $L$-包；并且应限制到在 [Lan79] 意义下不「反常」（anomalous）的自守表示（否则须把 $\Gal_F$ 换成 $\Gal_F\times\mathrm{SL}_2$，在 Arthur 关于非缓增内窥的猜想框架中工作 [Art89]）。关于该猜想更技术性的讨论，参看 [BG14]；更详尽的历史综述，参看 [Eme21]。

我们指出互反律与函子性之间的一个联系：由于函子性对（相容系统意义下的）伽罗瓦表示是平凡成立的，互反律的诸情形可用来推出函子性的诸情形。又由于 L-群与 C-群的定义中涉及伽罗瓦群，伽罗瓦群到代数群的表示有时可与 L-群的 L-同态或 C-群的 C-同态相关联，于是互反律的某些情形可并入函子性。（这正是 Langlands 对 Artin 猜想的原始视角，见 [Lan70] 的说明。）若把伽罗瓦群 $\Gal_F$ 换成假设中的 Langlands 群 $L_F$，把 $\Gal_F$-表示的相容系统换成 $L_F$ 的表示，则还可把互反律猜想推广到非代数的自守表示；或者（利用 L-群或 C-群的「$L_F$-形式」）把互反律完全并入函子性。（这后一视角的某个版本，Langlands 本人已时而采用，即使用 L-群的「Weil 群形式」。）

\ysrc{170}
\subsection{$p$-进视角}

本讲义并不打算按上述方式聚焦于互反律（即自守表示与动机、或与伽罗瓦表示的相容系统之间的关系），而是要固定一个素数 $p$，考察自守表示与 $p$-进伽罗瓦表示（对这样固定的 $p$）之间的关系。初看之下，这并不太改变问题：因为（至少对取值于 $\mathrm{GL}_d$ 的表示而言）半单伽罗瓦表示的相容系统由其中任一成员决定，而 Fontaine--Mazur 猜想给出了一个 $p$-进伽罗瓦表示成为动机的纯伽罗瓦论条件。

然而，聚焦于某个特定素数 $p$，会使自守形式与伽罗瓦表示理论中某些在更具动机色彩的互反律视角下不出现的侧面凸显出来。例如，至少自 Ramanujan 的工作以来，人们就知道自守形式可以满足模 $p$ 的幂的有趣同余。自 Swinnerton--Dyer [SD73] 与 Serre [Ser73] 的工作以来，人们已理解这些同余与 $p$-进伽罗瓦表示之间的类似同余相关（或者，若愿意，是其表现）。把自守形式之间同余的概念加以推广，人们便被引向取自守形式空间的各种 $p$-进完备化，从而得到 $p$-进自守形式的概念，其相伴的伽罗瓦表示不必是动机。与此相关，还有模 $p$ 自守形式的概念，可以给它配上模 $p$ 伽罗瓦表示；互反律猜想也已在此背景下被表述，著名的有 Serre [Ser87] 在经典模形式情形的工作，以及晚近他人更一般的工作（例如 [BDJ10]）。

我们完全无意综述这些发展；这里的目的只是指出：$p$-进视角是在很长一段时间里自然浮现的，并引出了一批自然的问题与关切。例如：如何得到某种自守形式的概念，使它容许 $p$-进整系数或模 $p$ 系数，并容纳真正的 $p$-进对象？又如何表述互反律，使之容纳 $p$-进整或模 $p$ 的伽罗瓦表示——这些表示将与自守侧引入的 $p$-进整或模 $p$ 对象相配？

关于 $p$-进理论，眼下再作最后一点说明：一旦允许对自守形式作 $p$-进插值，就会发现自守形式自然地处于族中（例如 Hida 族、Coleman 族……；综述见 [Eme11b]）。伽罗瓦表示同样处于族中（例如通过 Mazur 的形变理论 [Maz89]，他发展该理论至少部分是为回应 Hida 关于 $p$-进模形式常族的理论）。这些对象连续族的现象——在自守侧与伽罗瓦侧皆是——并不是互反律在其更具动机色彩的表述中的特征。（尖点自守表示是刚性对象；动机亦然，只要固定了某个数域作为定义域。）在转入下一话题之前，我们指出：伽罗瓦表示的族是几何 Langlands 纲领的一个特征。下文将回到这一点。

\subsection{互反律的证明}

关于互反律的一些最早结果是在像可解的情形下证明的，办法是把类域论（它在阿贝尔情形给出互反律）与更具经典自守色彩的技术结合起来（尤见 Langlands 与 Tunnell 的工作\ysrc{171}，[Lan80, Tun81]）。1995 年，Wiles [Wil95] 与 Taylor--Wiles [TW95] 证明了 $\Q$ 上半稳定椭圆曲线的模性，使我们对互反律的理解取得突破；这一结果此后不久即推广到 $\Q$ 上所有椭圆曲线 [BCDT01]。在此方法的基础上，互反律此后在许多其他有趣的情形中得到证明；关于过去 25 年的一些亮点，见 Calegari 的近期综述 [Cal23]。

[Wil95, TW95] 的方法以及此后建立在其思想之上的许多结果，其关键在于证明一个 $p$-进伽罗瓦表示的模性——或更一般地，自守性。更确切地说，Taylor--Wiles 方法给出了一条途径：由单个 $p$-进伽罗瓦表示 $\rho_1$ 的模性，推出另一个与 $\rho_1$ 模 $p$ 同余的 $p$-进伽罗瓦表示 $\rho_2$ 的模性；这通常称为模性提升（modularity lifting）。（我们从 $\rho_1$ 的模 $p$ 约化的模性——由于它与 $\rho_2$ 的模 $p$ 约化相同，这一模性承自 $\rho_2$——「提升」到 $p$-进表示 $\rho_1$ 本身的模性。）事实证明伽罗瓦表示之间存在大量这样的同余，模性可以由此从单个表示传播到许多其他表示（或用固定的素数 $p$，或用若干个不同的素数 $p$）。当然，必须先有一批模性已知的表示作为起点，例如 CM 形式，或像可解的二维 Artin 表示（后者用于 Wiles 对 Fermat 大定理的证明 [Wil95]）。

模性提升定理也称为「$R=\mathbf T$」定理，因为其证明方式是：把一个参数化与固定表示 $\rho_1$ 同余的伽罗瓦表示的 $\Z_p$-代数 $R$（一个「伽罗瓦形变环」）与另一个 $\Z_p$-代数 $\mathbf T$（一个「Hecke 代数」，即适当的（$p$-进）模形式空间的自同态代数）等同起来。过去 25 年间，人们逐渐认识到这类定理应在极大的一般性下成立，尽管要为 $R$ 与 $\mathbf T$ 找到恰当的定义、使字面上确有 $R=\mathbf T$，仍多少是一门技艺。特别地，人们已清楚（尤其继 Skinner--Wiles [SW99]、Wake--Wang-Erickson [WWE17] 与 Newton--Thorne [NT23] 之后）：一般情形下 $R$ 应取所谓的伪形变环（pseudodeformation ring），即伪表示的形变环，而非字面意义下伽罗瓦表示的形变环。

粗略地说，伪表示捕捉的是表示的特征多项式所给出的信息；因此它们也称为伪特征标。它们最初由 Wiles [Wil88] 对 $\mathrm{GL}_2$ 引入，后经许多作者对 $\mathrm{GL}_d$ 大为发展（我们特别指出 Chenevier 的行列式概念 [Che14]，它说明了如何对任意素数 $p$ 定义它们）。适用于一般约化群的理论由 Vincent Lafforgue [Laf18] 引入。

在 Lafforgue 看来（这一观点最近由 Fargues--Scholze [FS24, §VIII] 在局部伽罗瓦表示的设定中精确化），伪形变环 $R$ 是 L-参数模叠上整体函数之环，于是伪表示的模空间是 L-参数模叠的一个粗模空间；\footnote{此处及下文，「粗模空间」一词按通常的非正式方式使用；由于 L-参数的叠不是 Deligne--Mumford 叠，它们没有技术意义上的粗模空间，而只有 [Alp14] 意义下的适当模空间（adequate moduli space）。} 而对一般约化群，Hecke 代数 $\mathbf T$\ysrc{172}被替换为更大的代数，即「游历算子」代数（对 $\mathrm{GL}_d$ 它与通常的 Hecke 代数一致）。

从这个视角看，自然会问：能否把设想的等式 $R=\mathbf T$ 提升为关于 L-参数的实际模叠（而非其粗模空间）的陈述？例如，可以问 $p$-进模形式的空间是否容许用 L-参数叠上的层来解释。这类陈述的可能性及其局部类比，是本讲义的主题；但在进一步讨论这一可能性之前，我们先转向几何 Langlands 纲领——它也是 Vincent Lafforgue 工作的重要动机。

\subsection{非分歧几何 Langlands 与 Fargues--Scholze 猜想}

极其粗略地说，对连通约化群 $G$ 与曲线 $X$，非分歧几何 Langlands 对应取范畴等价的形式
\begin{equation}\tag{1.4.1}\label{eq:a4-1-4-1}
D(\mathrm{const.}\,\Bun_G(X))\cong \Qcoh(\LocSys_{\hatG}(X)),
\end{equation}
其中左端是 $X$ 上 $G$-丛的模叠 $\Bun_G$ 上「可构造层」的导出范畴，右端是 $X$ 上「$\hatG$-局部系统」的某个叠上「拟凝聚层」的导出范畴；例如见 [Gai17]。

Fargues--Scholze 的近期工作 [FS24] 通过「取 $X$ 为 Fargues--Fontaine 曲线」把经典局部 Langlands 对应几何化。（这样做的动机之一是：Fargues--Fontaine 曲线的基本群是 $\Gal_F$。）稍精确一些（但仍只是近似地）：若 $F/\Q_p$ 是有限扩张，$X$ 为 $F$ 对应的 Fargues--Fontaine 曲线，则对每个素数 $\ell\neq p$，人们预期有近似形式 (1.4.1) 的等价，其中左端换成 $\Bun_G$ 上（固体）$\ell$-进层的适当导出范畴，右端换成 $\Gal_F$ 的 $\hatG$-值 $\ell$-进表示之模叠上（ind-）凝聚层的导出范畴。

Fargues--Fontaine 曲线上 $G$-丛的叠 $\Bun_G$ 容许一个由局部闭子叠构成的分层，其中一个开层片（stratum）是商叠 $[\,\cdot\,/G(F)]$。该商叠上的 $\ell$-进层对应于 $G(F)$ 的光滑 $\ell$-进表示，故通过考虑从这个商叠到 $\Bun_G$ 的适当推前，(1.4.1) 形式的等价特别蕴含存在全忠实函子
\begin{equation}\tag{1.4.2}\label{eq:a4-1-4-2}
D(\mathrm{sm.}\,G(F))\to\Qcoh(\LocSys_{\hatG}),
\end{equation}
其中左端是 $G(F)$ 的 $\ell$-进表示的导出范畴，右端是 $\Gal_F$ 的 $\hatG$-值 $\ell$-进表示之叠上拟凝聚层的导出范畴。

特化到 $G=\mathrm{GL}_d$ 的情形，我们有下述粗略猜想，这里允许 $\ell=p$。

\begin{conj}{1.4.3}\label{conj:a4-1-4-3}
设 $F$ 是 $\Q_p$ 的有限扩张，$\mathcal O$ 是 $\Q_\ell$ 的某个有限扩张中的整数环（$p$ 与 $\ell$ 各表示某个固定素数），$d$ 是正整数。则 $\mathrm{GL}_d(F)$ 在挠 $\mathcal O$-模上的光滑表示范畴，容许全忠实嵌入到 $\mathcal X$ 上拟凝聚层的范畴中，其中 $\mathcal X$ 是参数化绝对伽罗瓦群 $\Gal_F$ 的 $d$ 维 $\ell$-进表示的适当模叠。
\end{conj}

\ysrc{173}我们之所以称此猜想为「粗略」，有几个理由。其一，须明确猜想陈述中出现的叠 $\mathcal X$：当 $\ell\neq p$ 时，应理解为 [Zhu25a] 引入的 Weil--Deligne 表示之叠（底层经典叠的其他构造亦见 [DHKM25] 与 [FS24]）；而在 $\ell=p$ 的情形，它应是 [EG23] 构造的平展 $(\varphi,\Gamma)$-模模叠之一（且先验地应把这些叠视为导出叠——不过在 $\ell\neq p$ 的情形，[Zhu25a] 已证明所得叠实际上是经典的；在 $\ell=p$ 的情形，Min [Min24] 借助 Böckle--Iyengar--Paškūnas [BIP23a] 的结果也证明了同样的事）。其二，须明确所考虑的范畴之种类：设想的嵌入不会与源与靶上自然的阿贝尔范畴结构相容，故应在适当的导出范畴（或者，多半更好，稳定 $\infty$-范畴）中工作（此处「适当」也颇有讲究——例如应在 $\mathcal X$ 上 ind-凝聚复形的三角/稳定 $\infty$-范畴中工作，然后对表示的三角范畴也作类似修改，如 [Zhu25a, §4.1] 那样）。

我们还想要猜想的一个版本：把 $\mathrm{GL}_d(F)$ 的光滑 $p$-幂挠表示换成局部解析表示（即与超收敛 $p$-进自守形式、Coleman 族等理论相关的那类 $p$-进表示），并把叠 $\mathcal X$ 换成 Fargues--Fontaine 曲线上等变向量丛的叠（它也可解释为 $(\varphi,\Gamma)$-模的叠，此时系数在 Robba 环中；见下文定理 5.1.5）。由于目前连把这一情形纳入上述粗略猜想（同时处理 $\ell\neq p$ 与 $\ell=p$ 两情形）的非正式办法都没有，我们把 $p$-进局部解析情形的猜想表述留到正文。只提一点：此情形的猜想可视为 $\ell=p$ 情形下猜想 1.4.3 的 $p$-进极限（由挠 $\mathcal O$-模过渡到完备无挠 $\mathcal O$-模得到）的超收敛版本。

\begin{rem}{1.4.4}\label{rem:a4-1-4-4}
据我们所知，$\ell=p$ 的、沿猜想 1.4.3 路线的猜想，最早由 Michael Harris 提出（见 [Har16, Question 4.7]）。我们最早是从 Laurent Fargues 那里得知：(1.4.1) 形式的 Fargues--Scholze 猜想应有 $p$-进版本，时在 2016 年夏。Peter Scholze 于 2018 年向 T.G. 解释了若干相关猜想；2019 年在波恩的 Hausdorff 学校，M.E.、T.G.、E.H. 与 Scholze 讨论了猜想 6.1.15 某版本的粗略表述。
\end{rem}

\begin{rem}{1.4.5}\label{rem:a4-1-4-5}
我们借 Fargues--Scholze 的工作来为猜想 1.4.3 提供动机。事实上，在 $\ell\neq p$ 的情形，沿猜想 1.4.3 路线的各种猜想已由 Ben-Zvi--Chen--Helm--Nadler、Fargues--Scholze、E.H. 与 Zhu 各自独立提出。这类猜想的历史见 [FS24, Zhu25a, Hel23, BZCHN24] 诸文的引言。自然可以猜测 $\ell=p$ 情形也应有类似猜想，但精确表述该是什么则不那么清楚。我们在第 6 节提出两个精确猜想，一个在「Banach」情形（聚焦于挠 $\mathcal O$-模上的表示，以及极限情形下 Banach 空间表示中的格），一个在「解析」设定（聚焦于局部解析表示），并解释它们之间的若干关系。
\end{rem}

\ysrc{174}然而，两个猜想似乎互不蕴含，我们也不知道它们是否有共同的加细。

这两种情形对应于 $p$-进自守形式（或 $p$-进模形式）的两个经典概念。Serre 把 $p$-进模形式定义为带挠系数的 $p$-进模形式之极限（更确切地说：模形式 $q$-展开式模 $p$ 的幂的约化之 $p$-进极限）。这些 $p$-进模形式空间当然属于「Banach 情形」；而 Katz 引入的超收敛模形式子空间则属于「解析情形」。在本讲义的讨论中，$p$-进模形式空间被替换为模曲线（或更一般的志村簇塔）的完备上同调，如 §3.2.11 与 §9.3.1 那样；超收敛模形式空间则被替换为完备上同调背后的局部解析向量空间，后者用于构造特征簇，见 §9.6。

\begin{rem}{1.4.6}\label{rem:a4-1-4-6}
读者可能会问：为什么我们只寻求全忠实函子，而不像上述 Fargues--Scholze 的工作那样寻求范畴等价？希望确实存在这样的范畴等价，似乎是合理的（甚至在整体设定中亦然，如 [FS25, §7] 所猜测）。困难之一在于，连 $\Bun_G$ 上 $p$-进层的范畴都还难以定义。这些困难在 Lucas Mann 的近期工作 [Man22] 以及 Lucas Mann 与 Johannes Anschütz、Arthur-César le Bras 的合作工作 [ALBM24] 中得到处理；后者将在 Johannes Anschütz、Arthur-César le Bras、Juan Esteban Rodriguez Camargo 与 Peter Scholze 关于解析合冲化（analytic syntomification）的即将发表的工作中得到推广。这些工作应会在表述精确猜想时起重要作用；但无论如何，在撰写本文时文献中还没有这样的精确猜想（更不用说定理了——我们指出，据我们所知，并不预期 Fargues--Scholze 对谱作用的构造会以任何简单方式搬到 $p$-进设定）。相反，我们可以对 $\mathrm{GL}_n(F)$ 表示范畴的全忠实嵌入给出精确猜想，甚至在 $\mathrm{GL}_2(\Q_p)$ 的情形证明定理。此外，至少对 $\mathrm{GL}_2(\Q_p)$ 与 $\mathrm{GL}_2(\Q_{p^2})$（及其内型），我们给出了（ind-）凝聚层范畴的半正交分解，并预期它们对应于 $\Bun_G$ 上适当层范畴由 $\Bun_G$ 自身闭包关系诱导的半正交分解。（此方向的一些结果见 §7.5。）
\end{rem}

\begin{rem}{1.4.7}\label{rem:a4-1-4-7}
没有特别的理由要在猜想 1.4.3 的陈述中限于群 $\mathrm{GL}_d$，而不考虑 $F$ 上更一般约化群的对偶群、L-群或 C-群（如 [BZCHN24, Hel23, Zhu25a] 那样）。然而，由于我们关心 $\ell=p$ 的情形，我们将大体聚焦于 $\mathrm{GL}_d$，因为相关的「参数叠」只在此情形被构造出来。（由于作为直觉与动机重要来源的 Taylor--Wiles 拼接文献也聚焦于此情形，这也为本文聚焦于它提供了另一理由。）
\end{rem}

\ysrc{175}
\subsection{$\ell$-进与 $p$-进情形的若干差异}

虽然可以用某种统一的方式表述猜想 1.4.3 的 $\ell$-进与 $p$-进版本，但两情形之间有显著差异。从表示论的角度看，有差异并不奇怪。事实上，仅对 $\mathrm{GL}_2(F)$，$\ell\neq p$ 时不可约光滑模 $\ell$ 表示的分类已相当清楚，且可对任意 $F$ 统一表述（这些结果属于 Vignéras，见 [Vig89]）。相反，$\ell=p$ 时只有 $F=\Q_p$ 情形的分类（属于 Barthel--Livné 与 Breuil [BL94, Bre03a]）；而当 $F$ 是 $\Q_p$ 的非平凡非分歧扩张时，会出现许多复杂情况：例如存在不落在主级数中的无限族容许不可约表示（由 Breuil--Paškūnas [BP12] 构造），以及（绝对）不可约但非容许的表示（由 Le [Le19] 构造）。此外，超奇异不可约容许表示从不是有限呈现的（见 Schraen 的论文 [Sch15b] 与 Wu 的 [Wu21]；这里「有限呈现」用的是定义 E.3.3 的意义）。

这些差异与困难在对应的伽罗瓦侧有镜像。当 $\ell\neq p$ 时，伽罗瓦表示的模叠是相当温和的零维 ind-代数叠：其 ind-结构来自存在无穷多个连通分支（由伽罗瓦表示的导子控制），而每个连通分支是仿射概形被约化群取商。相应地，这些叠有许多整体函数，并有可用伪表示描述的粗模空间（至少远离少数坏素数时如此；这已由 Fargues--Scholze 证明 [FS24, Thm. VIII.3.6]）。

在 $\ell=p$ 的 Banach 情形，我们考虑的模叠是 [EG23] 中定义的那些。它们也是 ind-代数叠，但其维数随 $[F:\Q_p]$ 增长，且 ind-结构复杂得多：事实上，它们自然是形式代数叠，只有有限多个连通分支（仅来自对中心特征的一个条件）。此外，我们不预期它们有任何有趣的整体函数；并且在 $\mathrm{GL}_2(\Q_p)$（以及环面）之外，我们不预期它们有非平凡的相伴模空间。

在解析情形，模叠在本讲义中定义（更详细的讨论将出现在即将发表的论文 [HHS] 中），即 Fargues--Fontaine 曲线上等变向量丛的模叠（等价地，Robba 环上 $(\varphi,\Gamma)$-模的模叠）。就维数而言，解析情形与 Banach 情形相似：（预期的）维数随 $[F:\Q_p]$ 增长，并应与 Banach 情形形式叠的一般纤维的维数一致（事实上，猜想地，这个一般纤维应与解析模叠的一个开子叠密切相关）。不过，这些叠现在不应有 ind-结构，而应是刚性解析 Artin 叠。尽管这些叠上有略多几个有趣的整体函数，我们预期它们都来自 Hodge--Tate 权理论，因而整体函数代数仍然丢失了模叠所含的大部分信息。

\begin{rem}{1.5.1}\label{rem:a4-1-5-1}
$\ell=p$ 与 $\ell\neq p$ 两情形的另一差异是对域 $F$ 的依赖。当 $\ell\neq p$ 时，（一般是猜想性的）函子在所有 $F$ 上有定性一致的描述，且为 $F=\Q_p$ 建立猜想的情形似乎并不比为一般 $F$ 更容易。相反，若 $\ell=p$，则我们只知道如何在 Banach 情形为 $\mathrm{GL}_2(\Q_p)$ 构造函子，此时它与 $\ell\neq p$ 的情形有一些定性上的相似之处\ysrc{176}（例如它「不太导出」，唯一的导出现象与有限维表示有关；又如超奇异不可约容许表示大致对应于支撑在不可约伽罗瓦表示上的摩天层）。

但已经对 $\mathrm{GL}_2(\Q_{p^2})$，我们预期在导出层面考虑函子是必要的，此时超奇异不可约容许表示不再对应于摩天层，而应对应于 $H^{-1}$ 非凝聚的复形。（见 §7.7。）
\end{rem}

\begin{rem}{1.5.2}\label{rem:a4-1-5-2}
$\ell$-进与 $p$-进设定的差异，多少让人想起几何 Langlands 的 Betti 版本与 de Rham 版本之间的差异（前者例如见 [BZN18]，后者见 [Gai15]）；不过出于我们的无知，我们不确定这一类比该被多认真地对待。
\end{rem}

\subsection{来自 Taylor--Wiles 拼接的动机}

Fargues 曾强调，§1.4 所述的猜想并非来自试图把几何 Langlands 纲领「复制粘贴」到 Fargues--Fontaine 曲线的设定，而是来自对 $p$-进周期态射、志村簇与 Rapoport--Zink 空间的 $p$-进几何、以及经典局部 Langlands 对应（特别是 $L$-包与内窥现象）的研究。类似地，虽然或许可以通过某种「在 Fargues--Scholze 猜想中令 $\ell=p$」的适当程序得到本文所呈的猜想，我们却是经多年迂回的过程得到它们的，其动机主要来自 $p$-进自守形式理论与 Taylor--Wiles 拼接方法。

Taylor--Wiles 拼接方法除了在塑造我们的想法方面重要之外，还提示了猜想 1.4.3 中那些猜想性函子的若干重要性质，即使在 $\ell\neq p$ 的情形亦然。例如，拼接方法提示此类函子的「积分核」应是一个真正的层（而非层的复形），这一性质在 Ben-Zvi--Chen--Helm--Nadler 与 Zhu 的构造的最初版本中完全看不出来。（Taylor--Wiles 方法与 $\mathrm{GL}_2(\Q_p)$ 的经典 $p$-进局部 Langlands 对应之间也有强联系，见 [CEG+18]。）

简言之，猜想 1.4.3 与 Taylor--Wiles 方法之间的联系如下。$\mathrm{GL}_d(F)$ 光滑表示的一个典型例子是紧诱导 $c\text{-Ind}_{\mathrm{GL}_d(\mathcal O_F)}^{\mathrm{GL}_d(F)}V$，其中 $V$ 是 $\mathrm{GL}_d(\mathcal O_F)$ 在真正有限的 $\mathcal O$-模 $V$ 上的连续（即光滑）表示。人们设想，猜想 1.4.3 的嵌入会把这样的表示送到真正的凝聚层（而非其复形），记为 $\mathscr A(V)$；并且这些凝聚层将是 Taylor--Wiles--Kisin 拼接方法中出现的拼接模 $M_\infty(V)$ 在 $\mathcal X$ 上的「整体化」，意思是：若 $R_\infty$ 是拼接理论中通常如此记的环之一，$\Spf R_\infty\to\mathcal X$ 是自然的 versal 态射，则 $\mathscr A(V)$ 沿此映射的拉回与 $M_\infty(V)$ 一致。特别地，该猜想蕴含拼接模应是「纯局部的」（此处「局部」应理解为处于局部伽罗瓦论的境况，与拼接时为整体自守/伽罗瓦论所作的选择相对立）。\ysrc{177}一般而言，判定这种「纯局部性」是否成立，是自守形式算术理论中众所周知的一大未解问题。

我们指出，在 [CEG+16] 的基础上，拼接策略已在 [BHS17b] 中被推广到「解析情形」。用与上面大致相同的表述，可以把拼接模配给局部解析表示，特别是局部解析主级数表示，即形如 $\Ind_B^G(\delta)^{\mathrm{an}}$ 的表示。相应的凝聚层应是猜想性「解析」函子在此类表示上取值之拉回（并且它与特征簇上有限斜率超收敛 $p$-进自守形式层在 $\delta$ 上的纤维相关，其方式类似于通常的拼接模与经典自守形式模相关）。在这一解析背景下，拼接模的「纯局部性」问题同样成为理论中的一大未解问题。

我们也可以朝相反方向进行。也就是说，鉴于拼接模 $M_\infty(V)$ 是由志村簇的上同调（系数在与 $V$ 对应的局部系统中）构造的，自然会问：志村簇的上同调能否以某种方式用猜想 1.4.3 的（猜想性）函子来计算。答案应当是肯定的，我们在 §9.2 中说明这一点（M.E. 与 Zhu 在准备中的论文 [EZ] 将有更详尽的解释）。与此并行，§9.6 用解析情形的猜想函子，对特征簇上超收敛自守形式层（可用志村簇的完备上同调定义）给出精确计算。

\subsection{「从伽罗瓦到自守」的方向}

「拼接模是纯局部的」这一预期有一个变体，即希望：可以对每个 $n$ 维 $p$-进伽罗瓦表示 $\Gal_F$ 的表示，典范地配上 $\mathrm{GL}_n(F)$ 的一个容许 Banach 表示；并且这一配应在志村簇的上同调中实现。在志村曲线情形，这一预期已有大量证据支持，最新的例子是 Breuil、Herzig、Hu、Morra 与 Schraen 关于 GK 维数的引人注目的结果 [BHH+23, BHH+21]。

虽然我们的焦点主要是从光滑表示到伽罗瓦表示叠上的层的函子，我们的猜想也蕴含存在从层到 $p$-进群表示的伴随函子。特别地，把它们作用于「摩天楼」层（即有限长、支撑在特定伽罗瓦表示上的层），我们便能说明我们的猜想与上述希望相容。（例如见注 7.8.6。）

\subsection{我们略去的话题}

由于时间与无知，我们不讨论 $p$-进 Langlands 纲领中许多近期（乃至不那么近期）的发展，但我们仍预期它们与这里讨论的想法有重要联系。（在 2025 年初修订本讲义时，我们克制了大幅更新内容的诱惑；下面这份清单反映的是最初讲课时略去的话题清单；如今它本可以显著加长！）特别地，我们对 Drinfeld 上半平面及相关 $p$-进局部对称空间的上同调只字未提，这方面有 Colmez、Dospinescu、Le Bras、Niziol 与 Pan 的研究 [CDN20, CDN23, CDN21, DLB17, Pan17]\ysrc{178}（感兴趣的读者可参看 \href{https://mathoverflow.net/questions/432932/how-does-the-cohomology-of-the-lubin-tate-drinfeld-tower-fit-into-categorical-p}{mathoverflow.net/questions/432932} 中关于预期联系的说明\tn{底本此处 URL 止于「categorical-p」，看似在印刷换行处被截断，译文照录；疑完整地址应续以「-adic-langlands-program」之类。}）；对 Scholze 的 $p$-进 Jacquet--Langlands 函子 [Sch18] 也只字未提；对志村簇的凝聚上同调同样如此，特别是与 Boxer 与 Pilloni 的高阶 Hida 与 Coleman 理论 [Pil20, BP22, BP21]、以及与 Diamond 与 Sasaki 的几何 Serre 权猜想 [DS23] 的任何关联。

\subsection{本讲义导读}

本讲义涵盖许多话题，我们希望各节大体可以彼此独立地阅读，必要时回查定义。在第 2 节建立一些记号之后，第 3 节以略带个人色彩的方式介绍 Taylor--Wiles 方法，以其与 $\mathrm{GL}_2(\Q_p)$ 经典 $p$-进局部 Langlands 对应的联系及它在启发我们猜想中的作用为线索。第 4 节简要回顾 [EG23] 中 $(\varphi,\Gamma)$-模的模叠之定义与基本性质。篇幅较长的第 5 节探讨这些构造在 Robba 环上 $(\varphi,\Gamma)$-模的「解析」背景中的类比。

第 6 节陈述我们主要的（局部）$p$-进 Langlands 猜想，兼及 Banach 与解析两种设定，并解释它们与文献中其他话题（例如 Breuil--Mézard 猜想）的关系。第 7 节解释这些猜想的已知情况；除相对直接的 $\mathrm{GL}_1$ 情形外，我们说明 Andrea Dotto 与 M.E.、T.G. 即将发表的工作，它证明了 $\mathrm{GL}_2(\Q_p)$ 的 Banach 猜想，并说明对 $F/\Q_p$ 非分歧的 $\mathrm{GL}_2(F)$，Banach 猜想的一些证据。第 8 节简要回顾 $\ell\neq p$ 情形的一些结果与预期。这些材料用于第 9 节，那里解释关于志村簇（完备）上同调与特征簇的猜想。

附录回顾（$\infty$-）范畴论与 $p$-进解析群表示论的若干结果，并证明另一些我们找不到文献出处的结果。

\begin{rem}{1.9.1}\label{rem:a4-1-9-1}
范畴 $p$-进 Langlands 对应仍处于襁褓之中，（正如本讲义所显示的）我们并不清楚我们所做的许多猜想的最终形式应该是什么。要使理论起步，还需要相当多的技术机器。尽管我们力求陈述精确而正确，我们的目标一直是尽可能解释与例示现象，而非详尽展开所有这些机器。特别地，我们大体上把 $\infty$-范畴与导出代数几何当作黑箱处理。

理论的未来发展极可能依赖凝聚态数学，特别是固体环（以及模、层……）的理论。把一些定义与构造用这些语言表述，也很可能使之更为简洁。然而，由于我们自身的局限，我们大体上没有尝试这样做。
\end{rem}

\subsection{致谢}

我们要感谢 IHES Langlands 纲领暑期学校的组织者（Pierre-Henri Chaudouard、Wee Teck Gan、Tasho Kaletha 与 Yiannis Sakellaridis）与参加者，感谢他们给予我们撰写本讲义的机会与鼓励，以及对我们多次修订的耐心。

% 本片止于印刷页 178 中段，续片 ch-a4-02 从 PDF 192 页起
